\section*{Chapitre \ref{ch:pl:reference-data-and-symbology-services} --- Données de référence et services de symbologie}

\begin{solution}{exo:pl:reference-data-and-symbology-services:1}
$384 / 10\,493 = 3.7~\%$. Cela arrive aux dates où le ticker appartenait à une autre cotation ou à aucune~: avant un renommage qui l'y attribue,
après un renommage qui l'en retire, et avant sa réattribution (\cref{prop:pl:reference-data-and-symbology-services:today}).
\end{solution}

\begin{solution}{exo:pl:reference-data-and-symbology-services:2}
Le CUSIP est resté attaché au titre lors du changement de ticker, si bien qu'un historique indexé par lui reste continu là où un historique indexé par ticker
se rompt (ou rattache l'historique d'un autre émetteur). Il reste néanmoins un identifiant externe~: il change avec certaines opérations sur titres,
est émis par une seule agence de numérotation pour un seul marché, et ne couvre pas tous les instruments que négocie la société~; celle-ci
indexe ses données par son propre identifiant permanent et stocke le CUSIP, l'ISIN, le FIGI et les tickers comme attributs datés.
\end{solution}

\begin{solution}{exo:pl:reference-data-and-symbology-services:3}
1 au jour 385 (rien de connu), 1 au jour 386 (annoncé mais non confirmé), 4 à partir du jour 389 (confirmé), et 4 au jour 412.
\end{solution}

\begin{solution}{exo:pl:reference-data-and-symbology-services:4}
La taille de lot mal saisie, 10, sur les six jours~: la règle de priorité recopie les erreurs de la source privilégiée. Le rapport de conflits la détecte, puisque
le fournisseur B indique 100 pour chacun de ces jours. Une règle de validation aurait bloqué l'erreur~: une taille de lot qui change d'un facteur dix
sans opération sur titres, ou qui n'est pas un lot standard du marché de la cotation, est mise en quarantaine pour examen et la \omterm{def:pl:reference-data-and-symbology-services:golden}{copie de référence} conserve
la valeur précédente entre-temps.
\end{solution}

\begin{solution}{exo:pl:reference-data-and-symbology-services:5}
Selon ce qui était connu au jour 255~: le fonds lui-même (l'acquisition n'était pas encore connue). Selon ce qui était connu au jour 262~: la société holding, par l'intermédiaire de la
banque. Le rapport reproduit doit montrer la première réponse, ce que la société savait lorsqu'elle a imprimé le rapport~; la seconde est une correction
à présenter séparément.
\end{solution}

\begin{solution}{exo:pl:reference-data-and-symbology-services:6}
Les splits non ajustés transforment chaque split en une baisse de $1 - 1/r$ dans la série de prix (un split quatre pour un ressemble à une baisse de $-75~\%$)~: sur les
bonnes cotations, le panier rapporte $-22.3~\%$ au lieu de $+23.5~\%$, soit 45.8 points d'écart. La correspondance d'aujourd'hui remplace cinq cotations visées par les
cotations qui portent aujourd'hui leurs tickers, avec leurs propres rendements depuis leur premier jour~: ajustés des splits, $+29.2~\%$, soit 5.7 points
au-dessus de la vérité. Ensemble, elles donnent $-16.6~\%$, soit 40.1 points en dessous~: les effets ne s'additionnent pas parce que les splits non ajustés tombent sur
les cotations visées, dont certaines ont été remplacées par la correspondance d'aujourd'hui au profit de cotations sans split.
\end{solution}

\begin{solution}{exo:pl:reference-data-and-symbology-services:7}
150 jours-instrument. Le fournisseur B se trompe sur 153, mais sur trois d'entre eux il n'a aucun ticker pour la nouvelle cotation, et la \omterm{def:pl:reference-data-and-symbology-services:golden}{copie de référence}
se replie sur le fournisseur A, qui a raison.
\end{solution}

\begin{solution}{exo:pl:reference-data-and-symbology-services:8}
Les rendements calculés à partir de prix ajustés sont corrects, mais un signal qui utilise les niveaux de prix ne l'est pas~: avec les facteurs d'aujourd'hui, le prix
d'avant un split est divisé par un ratio qui n'était pas connu à l'époque (pas même annoncé), si bien qu'une règle du type «~prix inférieur à 5~\$~» ou
«~proche d'un nombre rond~» voit des prix que le marché n'a jamais affichés. Il faut ajuster avec les facteurs connus à chaque instant de décision, ou
se limiter aux rendements.
\end{solution}

\begin{solution}{pb:pl:reference-data-and-symbology-services:1}
\begin{enumerate}
\item 12 renommages, 5 réattributions, 20 splits, dont un annulé.
\item La date de validité (à partir de laquelle le fait est vrai) et le temps de connaissance (quand la société l'a appris)~: la première sert à interroger
le passé, la seconde à reproduire ce qui était alors tenu pour vrai.
\item Deux renommages saisis avec deux jours de retard, et un renommage saisi seulement trois jours après que l'ancien ticker a été réattribué~:
146 jours avec un ticker erroné pour une seule cotation.
\item Une taille de lot de 10 au lieu de 100 pendant six jours.
\item Les valeurs des fournisseurs servent de preuves pour résoudre les conflits, pour mesurer chaque source et pour reconstruire la \omterm{def:pl:reference-data-and-symbology-services:golden}{copie de référence}
selon une nouvelle règle.
\item Jamais pour le fournisseur A, sur 153 jours-instrument pour le fournisseur B, jamais pour la \omterm{def:pl:reference-data-and-symbology-services:golden}{copie de référence}.
\item Le fournisseur A est prioritaire pour les tailles de lot et s'est trompé~; la \omterm{def:pl:reference-data-and-symbology-services:golden}{copie de référence} reprend sa valeur.
\item 150 pour les tickers, 6 pour les tailles de lot.
\item La cotation renommée au jour 33 (fournisseur A~: son nouveau ticker~; fournisseur B~: toujours l'ancien), le jour où son ancien ticker a été attribué
à une nouvelle cotation.
\item Garder A en premier pour les tickers (il ne s'est jamais trompé sur eux), et ajouter une validation pour les tailles de lot plutôt que de modifier la priorité~: A
a raison sur les tailles de lot partout sauf sur ces six jours.
\item 3.7~\%.
\item Cotation 56 jusqu'au jour 32, aucune cotation du jour 33 au jour 175, cotation 100 à partir du jour 176~; le fournisseur B la maintient sur la cotation 56 jusqu'à
le jour 178.
\item 1 avant la confirmation du jour 389, 4 après.
\item Il n'entre jamais dans le facteur~; l'annonce et l'annulation restent dans son historique.
\item Le fonds lui-même, tel que connu au jour 255~; la société holding, telle que connue au jour 262.
\item \textbf{Résultat de référence.} Le panier de 20 tickers rend $+23.5~\%$ en point-in-time, égal à la vérité~; avec la correspondance d'aujourd'hui et des prix non ajustés, il rend
$-16.6~\%$, soit 40.1 points faux, dont $-45.8$ points pour les seuls splits non ajustés, et
$+5.7$ pour la seule correspondance d'aujourd'hui.
\item Les splits~: une journée à $-75~\%$ se voit tout de suite. La correspondance d'aujourd'hui est plus dangereuse~: les prix des mauvaises cotations sont réels et
plausibles.
\item La cotation de chaque symbole à chaque date de décision et les facteurs d'ajustement, tous deux tels que connus à l'instant de décision.
\item Résoudre le cas dès le premier jour à partir d'un avis de la bourse, enregistrer la résolution et demander au fournisseur B d'améliorer
ses performances sur les renommages.
\item Parce que la société doit répondre à la fois à ce qui était vrai à une date et à ce qu'elle savait quand elle a agi, et que les deux diffèrent dès qu'un fait arrive
en retard ou est corrigé.
\end{enumerate}
\end{solution}

\begin{solution}{iq:pl:reference-data-and-symbology-services:1}
Les tickers changent et sont réattribués~: un historique indexé par ticker éclate le passé d'une même société entre plusieurs symboles et réunit des sociétés différentes
sous un même symbole. Il faut indexer par un identifiant permanent et conserver les tickers comme attributs datés.
\iqlookfor{Réattributions et renommages~; une clé permanente.}
\end{solution}

\begin{solution}{iq:pl:reference-data-and-symbology-services:2}
Par une priorité explicite pour ce champ, mesurée sur la fiabilité de chaque fournisseur, complétée par une validation~; signaler le conflit aux opérations et
le résoudre sur preuves. Conserver les deux valeurs des fournisseurs, la provenance de la valeur retenue et toute surcharge, de façon bitemporelle.
\iqlookfor{Priorité assortie de la provenance, conflits signalés, rien n'est écarté.}
\end{solution}

\begin{solution}{iq:pl:reference-data-and-symbology-services:3}
Ajuster les prix passés avec des facteurs connus seulement plus tard (un split annoncé après la décision) modifie les niveaux de prix que la stratégie a vus,
et une opération sur titres appliquée à sa date d'effet, alors qu'elle n'est connue que plus tard, antidate la connaissance.
Il faut stocker les opérations sur titres comme événements versionnés et demander les facteurs tels que connus à chaque décision.
\iqlookfor{Temps de connaissance de l'ajustement.}
\end{solution}

\begin{solution}{iq:pl:reference-data-and-symbology-services:4}
Les limites sont fixées sur les sociétés mères ultimes. Si une contrepartie est rachetée, l'exposition se regroupe sous une nouvelle mère à partir de la date
d'acquisition, mais la société ne l'a appris que plus tard~: un contrôle de limite effectué entre-temps était juste au vu de ce qui était connu, et doit pouvoir être reproduit ainsi,
tandis que la correction montre l'exposition réellement détenue par la société.
\iqlookfor{Validité confrontée au temps de connaissance, appliquée à une agrégation.}
\end{solution}

\begin{solution}{iq:pl:reference-data-and-symbology-services:5}
Sources~: fichiers de référence des bourses, deux fournisseurs ou plus, un flux d'opérations sur titres, une source sur les entités juridiques. Identifiants internes permanents
avec chaque identifiant externe (tickers par plateforme, ISIN, CUSIP, FIGI, codes des fournisseurs) comme attributs datés~; stockage bitemporel~;
une \omterm{def:pl:reference-data-and-symbology-services:golden}{copie de référence} par champ avec priorité, validation, surcharges et provenance~; une file d'attente des conflits~; les opérations sur titres comme événements versionnés.
Les consommateurs interrogent un service par identifiant ou symbole avec une date de validité et un temps de connaissance, jamais un fichier copié.
\iqlookfor{Identifiants, bitemporalité, \omterm{def:pl:reference-data-and-symbology-services:golden}{copie de référence} avec provenance, une interface de requête.}
\end{solution}
